\chapter{Контекст и мотивация}
\label{ch:motivation}

\section{Зачем нужна кибериммунная разработка, её место в горнодобывающей промышленности}
\index{кибериммунная разработка}

Современные киберфизические системы (промышленная автоматика, робототехника, транспорт) должны сохранять предсказуемое и безопасное поведение при ошибках, сбоях связи и целенаправленных воздействиях. Подходы \textbf{конструктивной информационной безопасности} и \textbf{кибериммунной архитектуры} требуют явного моделирования границ доверия, целей и предположений безопасности (ЦПБ), а также проверяемого покрытия критичного кода тестами.

Олимпиада в УГГУ направлена на то, чтобы участники на практике прошли путь от учебного прототипа до осмысленной архитектуры автономной буровой установки (АБУ) как системы с конструктивной информационной безопасностью (СКИБ) в составе сценария «\textbf{цифровой рудник}» (ЦР). 

Система \textbf{Регулятор} сертифицирует поставку программного обеспечения, при этом стоимость сертификации зависит от целей безопасности, объёма и сложности обеспечивающего эти цели безопасности кода. Кроме того, стоимость эксплуатации АБУ для ЦР также зависит от этих параметров.  

\section{Нормативная и методическая опора}
\index{ГОСТ Р 72118-2025}
\index{ISO/SAE 21434}

В учебном задании используются положения \textbf{ГОСТ~Р~72118-2025} (цели и предположения безопасности, доверенная вычислительная база, SBOM) и практики стандарта \textbf{ISO/SAE~21434} (моделирование угроз, TARA).

\section{Роли систем в учебном стенде}
\label{sec:roles}
\index{цифровой рудник}
\index{АБУ}
\index{Регулятор}

\begin{tabularx}{\textwidth}{@{}l X@{}}
\toprule
\textbf{Роль} & \textbf{Назначение в стенде} \\
\midrule
Цифровой рудник (ЦР) & Регистрация установок, миссии, политика сертификатов; обмен с АБУ по REST. \\
АБУ & Миссии, телеметрия; в заготовке --- технический долг (ДВБ не выделена). \\
Регулятор & Пакет сертификации (код, SGA, SBOM), тесты, покрытие, результат и идентификатор сертификата. \\
\bottomrule
\end{tabularx}

Связь между ролями и обмен данными подробнее раскрывается в главе~\ref{ch:architecture} и иллюстрируется диаграммой контекста (рис.~\ref{fig:context}).

\section{АБУ как СКИБ: ЦПБ ЦР, Регулятор и эксплуатант}
\label{sec:skib-why}
\index{СКИБ}
\index{доверенная вычислительная база}

\textbf{АБУ как СКИБ.} Автономная буровая установка выполняет критичные для безопасности и производства функции в контуре «рудник\,\textemdash\,установка». Проектирование АБУ как \textbf{системы с конструктивной информационной безопасностью} (СКИБ) означает, что архитектура изначально задаёт минимальные привилегии, явные границы доверия, выделение \textbf{доверенной вычислительной базы} (ДВБ) и проверяемое соответствие \textbf{целям и предположениям безопасности} (ЦПБ). Без такого проектирования добавление «умных» функций и зависимостей неизбежно размывает ответственность за поведение при сбоях и атаках: растёт поверхность угроз, а эксплуатант теряет предсказуемость.

\textbf{Цифровой рудник и цели безопасности.} ЦР --- не «прозрачный канал» между участниками, а надсистема с собственной политикой: регистрация установок, учёт миссий, правила допуска по сертификатам. Поэтому у ЦР тоже формулируются \textbf{цели безопасности} (ЦБ в идентификаторах SG\_\ldots{}; согласованные с ними ЦПБ): целостность и подлинность данных о допущенных к работе установках, управляемость сценариев миссий, согласованность с политикой сертификатов. Иначе компрометация или ошибка на стороне ЦР подрывает доверие ко всем подключённым АБУ, даже если каждая установка спроектирована корректно.

\textbf{Роль Регулятора.} \textbf{Регулятор} выступает независимым звеном оценки поставляемого пакета: код, SBOM, SGA, тесты и покрытие. Результат --- сертификат и измеримые затраты на сертификацию. Это создаёт для разработчика прозрачную обратную связь: усложнение ДВБ и зависимостей напрямую отражается в стоимости и времени допуска, а для эксплуатанта ЦР сертификат становится аргументом доверия к конкретной версии прошивки и составу компонентов.

\textbf{Значение для эксплуатанта ЦР и вектор развития ИИ.} Эксплуатанту важно не только «что умеет» установка, но и при каких предположениях о среде и надсистеме гарантируется безопасное поведение, что входит в ДВБ и что проверено тестами. Методология СКИБ (явные ЦПБ, TARA, трассируемость требований к тестам и к составу ПО) для автономных систем с элементами ИИ --- не формальность: с ростом автономности возрастают последствия ошибочных или скомпрометированных решений, и конструктивная ИБ задаёт язык, на котором эти риски можно сравнивать и снижать на этапе проектирования, а не постфактум.

\section{Основные задачи конкурсантов по шагам}
\label{sec:contestant-tasks}
\index{конкурсное задание}

Ниже --- ориентир по тому, \emph{что} делают участники на каждом этапе, \emph{что} намеренно не трогают без необходимости, \emph{как} проверяют результат и \emph{зачем} устроено именно так. Подробная процедура --- в главах~\ref{ch:certification}--\ref{ch:evaluation}.

\begin{itemize}
  \item \textbf{Изучение контекста и архитектуры.} Читают описание ЦР, АБУ и Регулятора, сценарии обмена. \textbf{Не меняют} код. \textbf{Проверяют} понимание вопросами к схеме и к регламенту. \textbf{Зачем:} без общей картины легко усилить угрозы там, где достаточно локального исправления, или наоборот --- сломать согласованность с надсистемой.

  \item \textbf{Развёртывание окружения.} Настраивают изолированное окружение Python (Pipenv) и зависимости из зафиксированного манифеста. \textbf{Не подменяют} системный интерпретатор пакетами «вручную». \textbf{Проверяют} воспроизводимость установки на чистой машине. \textbf{Зачем:} идентичность среды у жюри и у участника снижает споры о «у меня работает».

  \item \textbf{Запуск тестов.} Прогоняют полный набор тестов заготовки. \textbf{Не отключают} проверки без обоснования в отчёте. \textbf{Проверяют} зелёный прогон или фиксируют ожидаемые отличия. \textbf{Зачем:} тесты задают контракт ожидаемого поведения и регрессий при доработках.

  \item \textbf{Пакет сертификации и Регулятор.} Собирают пакет для сертификации. Проходят сценарий в учебном стенде. \textbf{Не игнорируют} SBOM и SGA как формальность. \textbf{Проверяют} прохождение цикла, идентификатор сертификата, соответствие заявленным ЦПБ. \textbf{Зачем:} цепочка «код\,---\,состав\,---\,цели\,---\,тесты\,---\,сертификат» моделирует реальный допуск к эксплуатации в составе ЦР.

  \item \textbf{Разработка решения.} Реализуют решение в выделенном каталоге исходников, опираясь на заготовку как на базовую линию. \textbf{Не переносят} произвольно код между зонами ответственности (например, смешивать контуры ДВБ и прочего без осознанного рефакторинга и обновления ЦПБ). \textbf{Проверяют} тестами, ручными сценариями миссии, повторной сертификацией при изменении поставки. \textbf{Зачем:} оценивается не «красота кода», а согласованность архитектуры, безопасности и доказательств.

  \item \textbf{Отчёт и сдача.} Заполняют отчёт по критериям, оформляют доступ к репозиторию. \textbf{Не подменяют} измеримые артефакты устными пояснениями без ссылок на код и логи. \textbf{Проверяют} полноту полей и воспроизводимость шагов для проверяющего. \textbf{Зачем:} жюри должно иметь возможность независимо верифицировать заявленные баллы.
\end{itemize}
